% chapters_parte2/05_mpls_qos.tex

\chapter{\eh{railway-track} MPLS y QoS}

El Internet público trata todos los paquetes igual: si hay congestión, se pierden
indiscriminadamente. Eso es inaceptable para voz, video en tiempo real, o los SLAs
de un cliente empresarial que paga por un enlace dedicado. \textbf{MPLS} permite
ingeniería de tráfico precisa; \textbf{QoS} garantiza que el tráfico crítico siempre
tenga recursos.

% ─────────────────────────────────────────────
\section{\eh{label} MPLS --- Multiprotocol Label Switching}

\begin{concepto}
En IP normal, cada router examina la dirección destino y hace una búsqueda en su
tabla de routing (\emph{longest prefix match}) para cada paquete. MPLS reemplaza
esa búsqueda repetida por un sistema de \textbf{etiquetas}: el primer router asigna
una etiqueta numérica al paquete y los routers intermedios solo hacen \emph{swap} de
etiquetas en O(1), sin buscar en tablas de routing completas.

\begin{itemize}
    \item \textbf{LSR (Label Switch Router):} router en el núcleo MPLS que solo hace
          swap de etiquetas.
    \item \textbf{LER (Label Edge Router):} router en el borde que asigna (push) o
          elimina (pop) etiquetas.
    \item \textbf{LSP (Label Switched Path):} el camino predeterminado que sigue un
          flujo de tráfico a través de la red MPLS.
    \item \textbf{FEC (Forwarding Equivalence Class):} conjunto de paquetes que reciben
          el mismo tratamiento de forwarding (misma etiqueta, mismo LSP).
\end{itemize}
\end{concepto}

\subsection{Formato de la etiqueta MPLS}

\begin{lstlisting}
 31       20 19    16 15   8 7        0
+----------+---------+------+----------+
|  Label   |   Exp   |  S   |   TTL    |
| (20 bits)|  (3 b)  | (1b) |  (8 b)   |
+----------+---------+------+----------+
\end{lstlisting}

\begin{itemize}
    \item \textbf{Label (20 bits):} identificador del LSP. Hasta $2^{20} = 1{,}048{,}576$ etiquetas.
    \item \textbf{Exp (3 bits):} Traffic Class para QoS (se mapea a DSCP).
    \item \textbf{S (Stack bit):} 1 = última etiqueta del stack (MPLS permite etiquetas anidadas).
    \item \textbf{TTL (8 bits):} igual que el TTL de IP, evita loops.
\end{itemize}

La etiqueta MPLS se inserta entre el header Ethernet (L2) y el header IP (L3),
por eso se llama protocolo de ``capa 2.5''.

% ─────────────────────────────────────────────
\section{\eh{arrows-counterclockwise} Operaciones de etiquetas}

\begin{center}
\begin{tabular}{lll}
\toprule
\textbf{Operación} & \textbf{Quién la hace} & \textbf{Acción} \\
\midrule
\textbf{Push}  & LER de ingreso  & Agrega etiqueta al paquete IP \\
\textbf{Swap}  & LSR en el núcleo & Reemplaza etiqueta por la siguiente del LSP \\
\textbf{Pop}   & LER de egreso (o penúltimo hop) & Elimina etiqueta, entrega IP normal \\
\bottomrule
\end{tabular}
\end{center}

\begin{ejemplo}
\textbf{Paquete de 192.168.1.1 $\to$ 10.0.0.1 a través de red MPLS de ISP:}
\begin{enumerate}
    \item \textbf{LER-A (ingreso):} reconoce FEC ``destino 10.0.0.0/8'', hace \textbf{push}
          etiqueta 200. Paquete: [L2][MPLS:200][IP:10.0.0.1][\ldots]
    \item \textbf{LSR-1:} etiqueta 200 $\to$ tabla dice swap a 350 hacia LSR-2.
          Paquete: [L2][MPLS:350][IP:10.0.0.1][\ldots]
    \item \textbf{LSR-2:} etiqueta 350 $\to$ swap a 120 hacia LER-B.
    \item \textbf{LER-B (egreso):} hace \textbf{pop}, entrega paquete IP puro.
          Forwarding normal basado en IP hacia 10.0.0.1.
\end{enumerate}
Los LSRs intermedios nunca miraron la IP destino. Solo swapearon etiquetas.
\end{ejemplo}

% ─────────────────────────────────────────────
\section{\eh{wrench} MPLS Traffic Engineering (TE)}

\begin{concepto}
Con routing IP normal, el tráfico siempre toma el camino de menor costo (shortest path).
Si hay un enlace congestionado pero otro libre con mayor costo, IP lo ignora. MPLS-TE
permite \textbf{forzar} LSPs por caminos específicos, balancear carga por múltiples
rutas, y garantizar ancho de banda reservado para flujos críticos.

El protocolo de señalización para MPLS-TE es \textbf{RSVP-TE} (Resource Reservation
Protocol - Traffic Engineering), que reserva recursos a lo largo del LSP.
\end{concepto}

\begin{moderno}
\emoji{rocket} \textbf{Segment Routing (SR-MPLS):} la evolución moderna de MPLS-TE.
En lugar de señalización RSVP hop-a-hop, el router de ingreso codifica el camino
completo como un \emph{stack de etiquetas} (segment list). No hay estado de LSP en
los routers intermedios. Google, Facebook y los grandes ISPs migraron de RSVP-TE
a Segment Routing por su simplicidad operacional.
\end{moderno}

% ─────────────────────────────────────────────
\section{\eh{balance-scale} QoS --- Quality of Service}

\begin{concepto}
\textbf{QoS} es el conjunto de mecanismos que garantiza comportamiento predecible
para tráfico crítico en una red compartida. Sin QoS, todos los paquetes compiten
igual por ancho de banda y, en congestión, cualquier flujo puede perder paquetes.
\end{concepto}

\subsection{Los tres problemas que QoS resuelve}

\begin{cuidado}
\begin{itemize}
    \item \textbf{Bandwidth:} un flujo puede consumir todo el ancho de banda disponible
          (Netflix descargando vs. llamada VoIP en el mismo enlace).
    \item \textbf{Delay (latencia):} colas llenas aumentan la latencia. Voz tolera máximo
          150ms one-way; un backup FTP puede tolerar segundos.
    \item \textbf{Jitter (variación de latencia):} paquetes de voz deben llegar con
          intervalos regulares. Jitter alto causa distorsión de audio.
\end{itemize}
\end{cuidado}

\subsection{Modelo DiffServ y DSCP}

\begin{concepto}
\textbf{DiffServ (Differentiated Services)} es el modelo dominante de QoS en Internet
y redes empresariales. Cada paquete IP lleva un campo \textbf{DSCP (Differentiated
Services Code Point)} de 6 bits en el header IP (parte del campo ToS/Traffic Class).
Los routers aplican tratamiento diferenciado según ese valor, sin mantener estado
por flujo.
\end{concepto}

\begin{center}
\begin{tabular}{llll}
\toprule
\textbf{Clase} & \textbf{DSCP} & \textbf{Valor} & \textbf{Uso típico} \\
\midrule
EF (Expedited Forwarding)   & EF   & 46 (101110) & VoIP, videoconferencia \\
AF41 (Assured Forwarding)   & AF41 & 34 (100010) & Video interactivo \\
AF31                         & AF31 & 26 (011010) & Video streaming \\
AF21                         & AF21 & 18 (010010) & Transacciones críticas \\
AF11                         & AF11 & 10 (001010) & Datos de negocios \\
CS0 (Best Effort)            & BE   &  0 (000000) & Tráfico normal \\
\bottomrule
\end{tabular}
\end{center}

\subsection{Mecanismos de QoS}

\begin{center}
\begin{tikzpicture}[
    box/.style={draw=cyan!50!white, fill=darkcardgray, text=textlight,
                rectangle, rounded corners=4pt, minimum width=3cm,
                minimum height=1cm, font=\small\bfseries, align=center},
    arrow/.style={->, thick, draw=gray!60!white},
]
    \node[box] (class) at (0,0)    {Clasificación\\{\tiny DSCP, ACL, NBAR}};
    \node[box] (mark)  at (4,0)    {Marcado\\{\tiny Set DSCP}};
    \node[box] (queue) at (8,0)    {Colas\\{\tiny LLQ, WFQ, CBWFQ}};
    \node[box] (sched) at (8,-2)   {Scheduling\\{\tiny FIFO, WRR, SP}};
    \node[box] (shape) at (4,-2)   {Shaping/Policing\\{\tiny Token bucket}};
    \node[box] (drop)  at (0,-2)   {Drop Policy\\{\tiny WRED, tail drop}};

    \draw[arrow] (class) -- (mark);
    \draw[arrow] (mark)  -- (queue);
    \draw[arrow] (queue) -- (sched);
    \draw[arrow] (sched) -- (shape);
    \draw[arrow] (shape) -- (drop);
\end{tikzpicture}
\end{center}

\begin{itemize}
    \item \textbf{LLQ (Low Latency Queue):} cola de alta prioridad para VoIP/video.
          Se sirve primero, siempre. Garantiza latencia mínima.
    \item \textbf{CBWFQ (Class-Based WFQ):} asigna porcentajes de ancho de banda
          garantizado por clase. Si nadie más usa su porción, puede tomar el exceso.
    \item \textbf{WRED (Weighted Random Early Detection):} descarta paquetes de baja
          prioridad antes de que la cola se llene, evitando TCP global synchronization.
    \item \textbf{Token Bucket:} modelo para policing/shaping. Se ``acumulan fichas''
          a una tasa CIR (Committed Information Rate); cada paquete consume fichas.
          Sin fichas $\to$ paquete se descarta (policing) o se retrasa (shaping).
\end{itemize}

\begin{ejemplo}
\textbf{QoS en una empresa con 100 Mbps hacia Internet:}
\begin{itemize}
    \item VoIP (DSCP EF): LLQ, 10 Mbps garantizados, latencia $<$10ms.
    \item Videoconferencia (DSCP AF41): CBWFQ, 30 Mbps garantizados.
    \item CRM/ERP (DSCP AF21): CBWFQ, 25 Mbps garantizados.
    \item Navegación general (DSCP BE): 35 Mbps restantes, WRED activo.
\end{itemize}
Si las llamadas VoIP no usan sus 10 Mbps, el resto puede aprovecharlos.
Pero si VoIP necesita esos 10 Mbps, los tiene garantizados sin importar lo demás.
\end{ejemplo}

\begin{moderno}
\emoji{rocket} \textbf{QoS en la nube:} AWS, GCP y Azure ofrecen Enhanced Networking
(SR-IOV) para reducir latencia en instancias. Google Cloud usa \textbf{Andromeda},
su stack de virtualización de red que aplica QoS a nivel de hipervisor. En Kubernetes,
los \texttt{LimitRange} y \texttt{ResourceQuota} son la capa de QoS de aplicación;
el QoS de red real requiere CNI plugins como Cilium con bandwidth enforcement via eBPF.
\end{moderno}
